\section{平台集成：快速、幂等、可审计}

模型和原则只有进入真实 upload/message path 才能产生效果。本节把儿童安全能力当作高风险 platform service：它需要 API、batch 或 stream integration，需要重试与降级，需要严格 tenant isolation，还要把审核、处置和报告状态连接成可追踪的 state machine。低延迟很重要，但不能用 silent failure 换取表面可用性。

\subsection{Detection service 与平台 action 分离}

平台可在 upload 前、存储后、公开分发前或用户举报后触发检测，不同位置具有不同 latency budget 和 context。Detection service 返回 match/score 与版本信息；policy engine 结合产品表面、地区、账户历史和用户申诉规则决定 hold、limit、remove、review 或 no action。分离两者有助于模型升级与 policy 变更独立审计。

\lecturefigure{09-platform-integration.png}{平台集成需要 detection、policy、action/report 与 audit 四层状态。}{本地字幕 03:20--05:50、10:10--16:20；Thorn platform solutions；概念重绘。}

\begin{knowledgebox}{读图：重试必须幂等}
\term{idempotency}（幂等性）表示同一事件重复执行不会产生额外副作用。Upload event 应携带稳定 event ID；detection result 以 content/model/version 去重；action service 防止重复删除或重复通知；reporting workflow 防止同一法定报告因 network retry 被多次提交。每层都要记录 request、result、retry 与 final disposition，才能在事故后重放。
\end{knowledgebox}

\begin{importantbox}{一个可操作的状态机}
可将对象状态写为
\[
\texttt{received}\rightarrow\texttt{screened}\rightarrow
\texttt{queued}\rightarrow\texttt{reviewed}\rightarrow
\texttt{actioned}\rightarrow\texttt{reported/closed}.
\]
其中每条边由明确 actor 和 policy version 触发；失败进入 retry、dead-letter 或 manual escalation，而不是把 timeout 当作 benign。对象可有并行的 content state、account state 和 report state，不能用一个布尔字段表示全部生命周期。
\end{importantbox}

\subsection{Failure mode 与降级策略}

本节从正常路径转向故障：第三方检测不可用、queue backlog、model version 回滚、reference update 失败、audit log 延迟或 reporting endpoint timeout 都可能发生。平台应预先定义 fail-open、fail-closed、hold-for-review 与 rate-limited sampling 的适用边界，并让 on-call 能看到积压、年龄、严重度和恢复估计。

\begin{warningbox}{Fail-open 与 fail-closed 都可能伤害用户}
Fail-open 可能让高风险对象继续传播；fail-closed 可能大面积阻断合法用户或造成不可逆处置。更稳健的方案通常是按信号置信度、产品表面和潜在伤害分层：高置信 known match 可进入强 hold，低置信 predictive signal 进入受控队列，系统性 outage 则启用速率限制、延迟发布和人工升级。
\end{warningbox}

\teachervoice{课堂把平台称为 detection、removal 与 reporting 的第一责任主体，并没有把技术供应商描绘成执法机关。这个边界要求产品集成返回可解释信号，同时让客户平台保留自己的 policy、用户关系和法定义务。}

\subsection{本章小结}

儿童安全能力必须作为 durable workflow 集成，而不是一次同步 API 调用。Detection、policy、action、reporting 和 audit 分层后，系统才能在重试、回滚、降级与调查中保持一致证据。

